\begin{recipe}{Pepernoten}
\kicker[Bakken,Zoet,Nederlands,Sinterklaas]{Bakken · zoet · Nederlands · sinterklaas}
\index[register]{Pepernoten}
\index[register]{Speculaaskruid}

\meta{15 minuten + 20--25 minuten bakken \dvd Ca.\ 60--70 stuks \dvd Oven 175\,°C}

\lettrine{W}{at} we in de winkel vaak als pepernoot kopen is eigenlijk een
kruidnoot, maar in de volksmond noemt iedereen ze gewoon pepernoten. Met deze
zes ingrediënten maak je ze zelf.

\blockrule{Ingrediënten}
\begin{ingredients}
  \ing{250 g}{zelfrijzend bakmeel}
  \ing{100 g}{koude ongezouten roomboter, in blokjes (40\%)}
  \ing{125 g}{donkere basterdsuiker (50\%)}
  \ing{4 tl}{koek- en speculaaskruiden}\index[register]{Speculaaskruid}
  \ing{3 el}{melk}
  \ing{1 snufje}{zout}
\end{ingredients}

\blockrule{Bereiding}
\tip{Bewaar de afgekoelde pepernoten in een goed afgesloten trommel, dan
blijven ze tot een week goed.}
\begin{steps}
  \item Verwarm de oven voor op 175\,°C (boven- en onderwarmte). Bekleed een bakplaat
        met bakpapier.
  \item Meng het bakmeel, de basterdsuiker, de koek- en speculaaskruiden en het
        zout in een grote kom.
  \item Voeg de blokjes koude boter toe en kneed met je vingertoppen door het
        meelmengsel tot er een kruimelige structuur ontstaat.
  \item Voeg de melk lepel voor lepel toe en kneed snel tot een stevige,
        samenhangende deegbal.
  \item Draai er kleine balletjes van, ter grootte van een flinke knikker, en
        leg ze met wat ruimte ertussen op de bakplaat. Ze hoeven niet platgedrukt
        te worden, dat gebeurt vanzelf in de oven.
  \item Bak 20 tot 25 minuten in het midden van de oven, tot de pepernoten gaar
        en lichtbruin zijn.
  \item Laat ze volledig afkoelen op een rooster. Pas na het afkoelen worden ze
        lekker knapperig.
\end{steps}
\end{recipe}
